\begin{definition}\term{Выпуклое множество} (\emph{convex set}) \(D\subseteq\mathbb R^p\) вместе с любыми \(x,y\in D\) содержит весь отрезок между ними: \((1-t)x+ty\in D\) при \(0\le t\le1\).\end{definition}
\begin{definition}\term{Выпуклая функция} (\emph{convex function}) \(f:D\to\mathbb R\) на выпуклом множестве удовлетворяет
\[f((1-t)x+ty)\le(1-t)f(x)+tf(y)\qquad(x,y\in D,\ 0\le t\le1).\]
График на отрезке лежит не выше соединяющей его концы хорды.\end{definition}
\begin{definition}\term{Строго выпуклая функция} (\emph{strictly convex function}) удовлетворяет строгому неравенству в предыдущем определении при \(x\ne y\), \(0<t<1\).\end{definition}
\begin{example}Функция \(f(x)=x^2\) на \(\mathbb R\) строго выпукла, поскольку
\[(1-t)a^2+tb^2-((1-t)a+tb)^2=t(1-t)(a-b)^2>0\]
при \(a\ne b\), \(0<t<1\). Строгая выпуклость даёт не более одного минимизатора: если два разных минимума достигнуты, их середина имеет ещё меньшее значение. Существование минимума само по себе из строгой выпуклости не следует.\end{example}
